\chapter{The Space of Candidate Definitions}
\label{ch:candidate-space}

\begin{plainlanguage}
\textbf{In plain terms:} This chapter builds the simplest possible filter for candidate definitions: collect evidence from different sources and keep only what's consistent with all of it. It's a natural first move, marked here as provisional, since the next chapter proves it cannot work. Along the way, it also gives a precise name to the sources of evidence: a \emph{perspective}.
\end{plainlanguage}

The first attempt at an elimination mechanism is the most natural one available, and it is worth building carefully, in full, before the next chapter shows why it fails --- because the way it fails turns out to determine the shape of everything that replaces it.

Before any elimination mechanism can be built, the thing doing the eliminating needs a name. This book is titled around perspectival convergence, and everything that follows talks constantly about observers, cultures, institutions, developmental stages, and rival theoretical schools as sources of evidence --- yet none of these have so far been given a single formal home. That gap is closed now, briefly, so that the object the title refers to actually exists in the notation rather than remaining a suggestive word floating above it.

\begin{definition}[Perspective]
A perspective $P$ is any source capable of producing evidence that bears on which candidate definitions in $\Theta$ remain defensible --- an individual observer, an interview subject, a culture's folk theory, an institution's recorded practice, a developmental stage, a scientific research program, or a system's demonstrated capacity or failure. A perspective need not be a person; it need only be a source from which a constraint can, in principle, be extracted.
\end{definition}

\begin{definition}[Induced Constraint]
Each perspective $P_i$ induces a constraint on the candidate space via a map
\[
\Phi : \{\text{perspectives}\} \to \{\text{constraints on } \Theta\}, \qquad C_i = \Phi(P_i).
\]
\end{definition}

$\Phi$ is not assumed to be simple, uniform, or fully specified in advance --- extracting $C_i$ from a given $P_i$ is itself substantive interpretive work, and different chapters of this book amount to case studies in what $\Phi$ looks like for different kinds of perspective (an ethnographic record, an institutional history, a developmental dataset). What matters for the formalism is only that the output of $\Phi$, whatever the input, is a constraint of the kind the rest of this chapter and the next can operate on. With this in place, perspectives are no longer a word the book uses informally around the edges of its math --- they are the explicit input to the function that produces everything $\Theta_n$, $\mu_t$, and later chapters' case studies are built from.

\subsection*{Criteria for perspectivehood}

The definition above is broad enough to invite a fair question: is a microscope a perspective? An LLM? A culture? A three-year-old child? Left unanswered, the breadth of the definition risks becoming a license to call anything a perspective whenever doing so is convenient, which would make the entire elimination apparatus that follows unfalsifiable in the same way an uncontrolled subsystem split would be (a concern this book takes seriously enough to build an explicit anti-vacuity principle for, later, in Chapter~\ref{ch:development}). Four criteria, taken jointly, are offered to draw the line.

A perspective must be capable of \emph{generating constraints} --- it must produce, in principle, some output that bears on which candidates in $\Theta$ remain defensible, not merely exist as an object that could in theory be studied. A bare microscope generates no constraint by itself; a microscope used by a researcher to test a specific structural hypothesis, as part of Chapter~\ref{ch:resolving-fork}'s distinctive-features case, does. It must be \emph{distinguishable from other perspectives} --- there must be a principled way to say where one perspective ends and another begins, which is what makes pairwise independence (Chapter~\ref{ch:resolving-fork}) a meaningful quantity to compute rather than an arbitrary partition imposed after the fact. It must be \emph{capable of disagreement} --- a perspective that is constitutionally incapable of producing a constraint inconsistent with any candidate contributes nothing to the elimination dynamics and is, for the purposes of this book, inert; this rules out perspectives defined so narrowly that they can only ever confirm what is already assumed. And it must be \emph{independently describable} --- nameable and specifiable without reference to the candidate it happens to be evaluated against, so that $\varepsilon_i(\theta)$ and $\mathrm{Ind}(i,j)$ in Chapter~\ref{ch:resolving-fork} can be assessed without circularity.

By these criteria, a culture, an institution, a developmental stage, and a research program all qualify, provided each can be shown (not merely asserted) to satisfy all four --- which is exactly the burden Chapters~\ref{ch:room-for-river} through \ref{ch:recursion} attempt to discharge case by case, rather than assuming met by default. An LLM qualifies under the same test only to the extent that its outputs can be shown to constrain $\Theta$ independently of the training data and assumptions a human perspective already supplied --- a question this book does not resolve and flags explicitly as open, since the independence criterion is precisely what would need to be established, not assumed, before treating a model's output as a perspective in this technical sense rather than a restatement of perspectives already counted.

With Perspective and Induced Constraint in hand, the naive elimination mechanism can now be stated precisely.

\begin{naivemodel}
\begin{definition}[Naive Elimination]
Let $n$ perspectives $P_1, \ldots, P_n$ be consulted in sequence, each inducing a constraint $C_i = \Phi(P_i) \subseteq \Theta$ --- the subset of candidate definitions consistent with what that perspective has shown. Define
\[
\Theta_n = \bigcap_{i=1}^{n} C_i.
\]
\end{definition}

This is elimination in its most literal form. Each new perspective rules out whatever it is inconsistent with, and the survivors are whatever remains consistent with everything encountered so far. It has an obvious appeal: it requires no weighting, no judgment calls about which perspectives matter more than others, no machinery beyond set intersection. Every constraint counts equally; every constraint that rules something out, rules it out permanently.

\begin{proposition}[Monotonicity]
For pure intersection, $\Theta_{n+1} \subseteq \Theta_n$ for every $n$.
\end{proposition}

This follows immediately from the definition of intersection --- adding one more set to an intersection can only remove elements, never add them --- and the proof is not worth more than the sentence just given. It is stated explicitly because it is about to become the chapter's entire problem, and a problem is easiest to see once its source has been isolated and named. \emph{This model is proven unworkable in Chapter~\ref{ch:nonmonotonic} and repaired in Chapter~\ref{ch:fidelity}.}
\end{naivemodel}
